\documentclass[12pt]{article}
\usepackage{amsmath} % For advanced math typesetting
\usepackage{amsfonts} % For additional math fonts
\usepackage{amsthm} % For theorem environments
\usepackage{graphicx}

\newtheorem{definition}{Definition}

\title{EIA GDP Indirect Impact Analysis}
\author{Navit Dori}
\date{\today}

\begin{document}
\maketitle


\section{Basic Leontieff equations}

\begin{equation}
\begin{array}{ll}  % 'll' ensures both columns are left-aligned
    x   & \quad \text{Output} \\
   II   & \quad \text{Intermediate input, also called IO matrix} \\
    f    & \quad \text{Final demand} \\
    T   & \quad \text{Technical coefficient matrix, sometimes called A} 
\end{array}
\end{equation}
x is the total output vector, which is the total output productive sectors produce,
\[
\vec{x} =
\begin{pmatrix}
x_1 \\
x_2 \\
\vdots \\
x_n
\end{pmatrix}
\]

II is the intermediate input matrix, or the IO matrix, which is what productive sectors buy and sell from other productive sectors,

Each element in the intermediate input matrix $II$ is devided by the total output of the column (total output of the buying sector) to get the technical coeeficient:
\begin{equation}\label{eq:Tdef}
T_{ij} = \frac{II_{ij}}{x_j}
\end{equation}
This is an element-by-element multiplication, not a matrix multiplication. It is equivalent to
\begin{equation}
II_{ij} = T_{ij}x_j.
\end{equation}


The fundamental equation states that the total output is a sum of the intermediate input and the final demand,
\begin{equation}
\begin{aligned}
    x_i &= II_i+f_i \\
         &= \sum_k {T_{ik} x_k} + f_i, 
\end{aligned}
\end{equation}
where $II_i=\sum_k{T_{ik}x_k}$ is what sector $i$ sells to the other sectors, a sum of the row correxponding to sector $i$.  $f_i$ is what sector $i$ sells to households.

In matrix form the equation becomes
\begin{equation}
\vec{x} = T \vec{x} + f
\end{equation}


leading to
\begin{equation}
\vec{x} = (I-T)^{-1}f.
\end{equation}

We usually refer to $(I-T)^{-1}$ as the Leontieff matrix, $L$.

\section{Value Added}
GDP is the same as value added
\begin{equation}
\text{Value Added} = \text{Output} - \text{Intermediate Input} 
\end{equation}

\begin{equation}
\text{Value Added} = \text{Wages} + \text{Taxes} - \text{Subsidies} + \text{Rent} +\text{Capital} + \text{Profit}
\end{equation}

\section{GDP Direct Impact}

\begin{equation}
\begin{aligned}
\text{Direct Impact} &= \text{Wages} + \text{Taxes (of production)} - \text{Subsidies (of production)} \\
&\quad + \text{Operating Surplus} + \text{Capital} + \text{Profit},
\end{aligned}
\end{equation}

or
\begin{equation}
\text{GDP}_{\text{Direct Impact}} = \frac{VA}{Output} = \frac{GDP}{Output}
\end{equation}

\begin{equation}
VA = GDP
\end{equation}
%
The above formula has the advantage that for any new total output vector $\vec{x'}$ the new GDP is $GDP_{\text{Direct Impact}}\vec{x'}$, which is to say that for sector $i$ the new GDP is $GDP_{\text{Direct Impact},i}\vec{x'}_i$
%

\section{GDP indirect impact - our EIA model version}

Here I follow the power point file \textit{EIA Modelling Workshop-Day 1\_Dec 23 2019.pptx}. 
Reetika explained over the phone that the vector $E'$ is based on $II' = Outcome' - Value Added'$. 
All variable denoted by $'$ refer to variables after the shock. 
In the example given in the file (slide 21), 
\begin{equation}
II'_a = Output'_a - Value Added'_a = 5-3 = 2. 
\end{equation}
Value added is refered to as \textbf{GDP direct} in the file.

The next step is
\begin{equation}
E'_a = T_{i1}(output'_a - \text{GDP}'_a)  = T_{i1} \cdot \text{II}'_a
\end{equation}
where $T_{i1}$ is the technical coefficients of the first column of T, refering to sector $a$.
We could write $T_{i1} = T_{ia}$.
The first column of $T$ is $[0.00, 0.17, 0.33]$ and so we have
\begin{equation}
E' = [0.00, 0.17, 0.33] \cdot 2 = [0.00, 0.34, 0.66]
\end{equation}

The next step is to multiply $E'$ by the inverse Leontieff matrix $L$ as if it were a final demand variable. The equation given in slide 21 is
\begin{equation}
Y' = (I-A)^{-1} E'
\end{equation}
where in the file's notation $Y'$ is the total output (here $x'$) and $A$ is the technical coefficient (here $T$).

$Y'$ is thus
\begin{equation}
Y' = \begin{pmatrix}
1.154 & 0.385 & 0.269 \\
0.322 & 1.275 & 0.319 \\
0.503 & 0.385 & 1.325
\end{pmatrix}
\begin{pmatrix}
0.00 \\
0.34 \\
0.66
\end{pmatrix}
=
\begin{pmatrix}
0.31 \\
0.64 \\
1.01
\end{pmatrix}
\end{equation}

The file continues with
\begin{equation}\label{eq:typeI_Reetika}
\text{GDP} type I = Y'  \cdot \text{GDP}_{\text{ratio}}
\end{equation}

where $\text{GDP}_{\text{ratio}}$ is calulated in slide 18 to be
\begin{equation}
\text{GDP to output ratio of sector j}  = \frac{\text{GDP}_j}{x_j}
\end{equation}
and the result is
\begin{equation}
\text{GDP to output ratio}  = [0.50, 0.47, 0.55]
\end{equation}
However, slide 21 uses a different vector
\[
\text{GDP}ratio = [0.44, 0.37, 0.50]
\]

The final step is to carry out the multiplication of eq. \ref{eq:typeI_Reetika} and sum over the result
\begin{equation}
\begin{aligned}
\text{GDP} type I_a &= Y' \cdot \text{GDP}ratio \\
                    &= [0.31, 0.64, 1.01] \cdot [0.44, 0.37, 0.50] \\
                    &= [0.14, 0.24, 0.50]
\end{aligned}
\end{equation}
\begin{equation}
\text{GDP} indirect_a = \sum_i{ Y'_i \cdot \text{GDP}ratio_i }= 0.14+0.24+0.50 = 0.88
\end{equation}
 





%%%%%%%%%%%%%%%%%%%%%%%%%%%%%%
\end{document}
